\documentclass[11pt]{article}
\usepackage[margin=1in]{geometry}
\usepackage{amsmath,amssymb,amsthm,booktabs,microtype,hyperref}
\hypersetup{colorlinks=true,linkcolor=blue,urlcolor=blue,citecolor=blue,
 pdftitle={Growing powers of Mahonian coefficients},pdfsubject={OEIS A380274, A380275: growing-power theta crossover}}
\usepackage{array}\newcommand{\repo}[1]{\nolinkurl{#1}}
\newenvironment{writenote}{\begin{quote}\small\sloppy\textbf{[write, 2026-10-02]}\ }{\end{quote}}
\newtheorem{theorem}{Theorem}\newtheorem{lemma}[theorem]{Lemma}\newtheorem{corollary}[theorem]{Corollary}\newtheorem{remark}[theorem]{Remark}
\newcommand{\V}{V_n}\newcommand{\T}{T_{n,r}}\newcommand{\Z}{\mathbb Z}\newcommand{\E}{\mathbb E}\newcommand{\dd}{\,d}
\title{Growing powers of Mahonian coefficients}
\author{Research note}\date{1 October 2026}
\begin{document}\maketitle
\begin{abstract}
The sum of the $r$th powers of Mahonian coefficients changes from a Gaussian integral to a discrete Gaussian partition function when $r$ is comparable to the inversion variance, of order $n^3$. We prove the parity-sensitive transition, both adjacent growing-power regimes, a uniform arbitrary-order expansion on the transition scale, and an explicit fourth-order correction. A fixed critical path has a controlled Lambert-branch inverse. We also remove the exact modal normalization to obtain a multiplicative equivalent with quadratic and linear exponential corrections. These are growing-power statements; the fixed-power equivalent and three corrections already occur in Wang's 2026 manuscript. Priority for the growing-power extension is not established by the targeted literature search accompanying this note.
\end{abstract}

\section*{Report note (ProveIt, 2 October 2026)}\phantomsection\label{mgp:sec:note}
\begingroup\small\sloppy
\noindent\textbf{Source.} This report prints one manuscript in full: batch 77 of ProveIt's
\texttt{docs/incoming}, manuscript 55, archive \texttt{oeis-mahonian-crossover-result.zip} (main file
\texttt{mahonian\_crossover.tex}, 10-page PDF, dated 1 October 2026, author line ``Research note''),
arrival commit \texttt{096ee7b87}, placement \texttt{aa7345800}. It names no ProveIt commit and
cites no repository file. Whether it was AI-assisted is not stated. It is unrefereed and not
formalized; its shipped audit (\texttt{audit.md}) is the delivery's own review.

\smallskip\noindent\textbf{Changes.} Every label carries the prefix \texttt{mgp:}. This note and three
dated notes marked \textbf{[write]} (Sections 7, 8 and 9) were added, with PDF metadata. No
statement, proof, symbol or number was changed; theorem and section numbers are the manuscript's.

\smallskip\noindent\textbf{Priority.} The fixed-power asymptotics of $S_n(3)$ and $S_n(4)$ (OEIS
A380274, A380275), with three corrections, are Wang's \cite{wang}; this note claims only the
growing-power extension, and its targeted search does not establish priority for that either
(abstract and Section 9, kept verbatim).

\smallskip\noindent\textbf{Letters used twice.} The manuscript reuses some letters; the meaning is
always fixed by the section.
\begin{center}
\begin{tabular}{@{}>{$}l<{$}p{0.4\textwidth}p{0.4\textwidth}@{}}
\toprule
\text{Letter} & One meaning & Another meaning\\
\midrule
Y_n & Corollary 2: the common inversion count minus $\mu$ & Section 7: $\log S_n(\rho\V)$\\
q & $q_n=M_n(\lfloor\mu\rfloor-1)/C_n$, the first modal ratio & argument of $Z_\delta(q)$, $\Theta_\delta(q)$; the variable of the Mahonian product\\
V & $\V=\sigma_n^2$, the inversion variance & $V(z)$ (formal, Section 5) and $V(x)$ (real model, Section 7)\\
f,\ g & $f_n$, the Fourier interpolation & $f(x),g(x)$ in the proof of Theorem~\ref{mgp:thm:inverse}\\
J & $J_{h,n}$, exact weighted Fourier moments & $J_h(z)$, their formal series\\
A,\ B & $A(u,z)$, $A_j$; $B_{\delta,L}$ & arc and cutoff constants $A$; interval bound $B$; Bernoulli $B_{2m}$\\
H & $H(x,z)$, formal profile & R\'enyi entropies $H_r$, $H_\infty$ (Section 9)\\
N & $N=n(n-1)/2$, the maximal inversion count & (no threshold index is called $N$)\\
\bottomrule
\end{tabular}
\end{center}
The theta functions $Z_\delta$, $\Theta_\delta$ of \eqref{mgp:eq:theta} have the lattice shift as
subscript; the collection's ballot-sum report uses a theta function whose subscript is a scale (dated
note in Section 9).

\smallskip\noindent\textbf{Not claimed} (all stated in the text): any fixed-power result; convergence
of the infinite expansion, exponentially small remainders or a resurgent transseries
(Section 5); a parity-independent critical-window limit; a certified onset for the rounding of the
inverse, or a ceiling rule for arbitrary thresholds (Section 7); novelty of the Fourier-curvature
method, of discrete Gaussian theta normalizers or of high-order entropy limits (Section 9); priority.
The numerical tables check algebra and parity conventions; the error bounds come from the proofs.
\endgroup

\section{Definitions and results}
Write
\[
 M_n(k)=[q^k]\prod_{j=1}^n(1+q+\cdots+q^{j-1}),\qquad
 S_n(r)=\sum_{k=0}^{N}M_n(k)^r,\qquad N=\frac{n(n-1)}2.
\]
Here $r>0$ can be real; for integral $r$, the sum counts ordered $r$-tuples of permutations with the same inversion number. Set
\[
 \mu=N/2,\quad \V=\sigma_n^2=\frac{n(n-1)(2n+5)}{72},\quad
 \delta=\mu-\lfloor\mu\rfloor,\quad C_n=M_n(\lfloor\mu\rfloor).
\]
Thus $\delta=0$ for $n\equiv0,1\pmod4$ and $\delta=1/2$ for $n\equiv2,3\pmod4$. The exact normalized sum and two theta functions are
\begin{align}
 T_{n,r}&=S_n(r)/C_n^r, &\rho_n&=r/\V,\label{mgp:eq:T}\\
 Z_\delta(q)&=\sum_{x\in\Z+\delta}e^{-qx^2/2},&
 \Theta_\delta(q)&=e^{q\delta^2/2}Z_\delta(q).\label{mgp:eq:theta}
\end{align}
All asymptotic statements below have $n\to\infty$.

\begin{theorem}[Complete growing-power leading behavior]\label{mgp:thm:main}
For every sequence $r_n\to\infty$,
\[
 \frac{T_{n,r_n}}{\Theta_{\delta}(r_n/\V)}\longrightarrow1.
\]
In particular:
\begin{enumerate}
\item If $r_n/\V\to0$, then $T_{n,r_n}\sim\sqrt{2\pi\V/r_n}$.
\item If $r_n/\V\to\rho\in(0,\infty)$ along a fixed parity class, then $T_{n,r_n}\to\Theta_\delta(\rho)$.
\item If $r_n/\V\to\infty$, then $T_{n,r_n}\to1+2\delta$.
\end{enumerate}
More precisely, put
\[
 q_n=\frac{M_n(\lfloor\mu\rfloor-1)}{C_n}.
\]
In the third regime,
\[
 T_{n,r_n}-(1+2\delta)\sim2q_n^{r_n}.
\]
\end{theorem}
\begin{corollary}[Common inversion count in a collision]
For integral $r_n\to\infty$, condition $r_n$ independent uniform permutations of size $n$ to have the same inversion count, and let $Y_n$ be that common count minus $\mu$. In the critical window $\rho_n\to\rho\in(0,\infty)$ with fixed $\delta$, its distribution converges in total variation to
\[
 \mathbb P(Y=x)=\frac{e^{-\rho x^2/2}}{Z_\delta(\rho)},\qquad x\in\mathbb Z+\delta.
\]
In the subcritical regime of Theorem~\ref{mgp:thm:main}, $\sqrt{r_n}Y_n/\sigma_n$ converges in distribution to a standard normal. In the supercritical regime, the conditional distribution concentrates uniformly on the one or two modes.
\end{corollary}
The conditional masses are exactly $M_n(\mu+x)^r/S_n(r)$. The assertions follow from the same geometric majorant, Gaussian Riemann sum, and modal tail bound used in the proof of Theorem~\ref{mgp:thm:main}.

The one or two modes must be retained: there is no single parity-independent critical-window limit. Also, replacing $C_n$ by its leading normal approximation before taking the $r_n$th power loses an exponential factor of order $\exp(cn^2)$ in the critical window.

\begin{theorem}[An explicit uniform correction]\label{mgp:thm:fourth}
Define
\begin{align*}
 a_{[4]}(n)={}&1-\frac{27}{25n}+\frac{21087}{61250n^2}
 -\frac{19894023}{3062500n^3}
 +\frac{18889730796483}{907878125000n^4},\\
 \Psi_\delta(q)={}&\sum_{x\in\Z+\delta}(x^4-\delta^4)
             e^{-q(x^2-\delta^2)/2}.
\end{align*}
Uniformly for $\rho_n\in K$, where $K$ is any compact subset of $(0,\infty)$,
\[
 T_{n,r}=\Theta_\delta\bigl(\rho_n a_{[4]}(n)\bigr)
 -\frac{81\rho_n}{25n^4}\Psi_\delta(\rho_n)+O_K(n^{-5}).
\]
Here $\Psi_\delta=4\Theta_\delta''-4\delta^2\Theta_\delta'$. Arbitrary fixed orders, with uniform $O_K(n^{-L-1})$ remainders after order $L$, are generated in Section~\ref{mgp:sec:allorders}.
\end{theorem}
Consequently the first three inverse-$n$ orders of the critical-window expansion arise entirely from a curvature adjustment. The first nonquadratic local-profile contribution occurs at order $n^{-4}$.

\begin{theorem}[An unnormalized multiplicative equivalent]\label{mgp:thm:unnorm}
Uniformly for $\rho_n=r/\V$ in a compact subset of $(0,\infty)$,
\begin{align}
 S_n(r)={}&\left(\frac{n!}{\sqrt{2\pi\V}}\right)^r
 \exp\left\{-\rho_n\left(\frac{3n^2}{400}
                  +\frac{5883n}{490000}
                  +\frac{13963}{500000}\right)\right\}
 Z_\delta(\rho_n)\{1+O_K(n^{-1})\}.\label{mgp:eq:unnorm}
\end{align}
A relative expansion through any fixed order $L$ follows by retaining the logarithm of the central Fourier integral through order $n^{-L-3}$, together with the lattice-profile expansion through order $L$.
\end{theorem}
One must use the actual $\rho_n$ in the finite-$n$ exponential. Merely knowing $\rho_n\to\rho$ does not justify replacing it there by $\rho$; an $o(1)$ change can be magnified by $n^2$. Exact $n!$ in~\eqref{mgp:eq:unnorm} is elementary and can be replaced by a sufficiently deep Stirling expansion.

\section{Fourier estimates with microscopic cancellation}
The inversion count is distributed as $X_n=\sum_{j=1}^n U_j$, where the $U_j$ are independent and uniform on $\{0,\ldots,j-1\}$. The centered characteristic function is
\[
 \phi_n(t)=\E e^{it(X_n-\mu)}
 =\prod_{j=1}^n\frac{\sin(jt/2)}{j\sin(t/2)}.
\]
It is real and even. Define the real Fourier interpolation and its standardized version by
\[
 f_n(x)=\frac1{2\pi}\int_{-\pi}^{\pi}\phi_n(t)\cos(xt)\dd t,
 \qquad F_n(y)=\sigma_n f_n(\sigma_n y).
\]
For $x\in\Z+\delta$ within the support, $f_n(x)=M_n(\mu+x)/n!$. The interpolation need not be a probability density on the whole real line; only its integral representation and positivity in a bounded standardized window will be used.

\begin{lemma}[Fourier bounds and finite-order expansions]\label{mgp:lem:fourier}
There exist positive absolute constants $c,C$ such that, for all sufficiently large $n$,
\[
 |\phi_n(t)|\le
 \begin{cases}
 e^{-c\V t^2},& |t|\le1/n,\\
 Ce^{-cn},&1/n\le|t|\le\pi.
 \end{cases}
\]
Every fixed derivative of $F_n$ has a uniform finite-order expansion in $n^{-1}$ on bounded real intervals, with an error of the next order. In particular,
\[
 F_n(y)=\frac{e^{-y^2/2}}{\sqrt{2\pi}}+O_B(n^{-1})
 \quad (|y|\le B),
\]
with the same type of estimate after any fixed number of derivatives. Thus $F_n>0$ there for large $n$, and
\[
 \frac{\dd^2}{\dd y^2}\log F_n(y)=-1+O_B(n^{-1}).\label{mgp:eq:curv}
\]
\end{lemma}
\begin{proof}
For $|jt|\le1$, the logarithm of the absolute sine ratio is at most $-c(j^2-1)t^2$, for a fixed positive $c$; this follows either by the Taylor series of $\log(\sin z/z)$ or compactness after division by $z^2$. Summing proves the first bound. For $1/n\le|t|\le A/n$, retain only the factors with $j\le\lfloor1/|t|\rfloor$. Their product is at most $\exp(-c/|t|)\le\exp(-cn/A)$; the remaining factors have modulus at most one. For $A/n\le|t|\le\pi$, use
\[
 \left|\frac{\sin(jt/2)}{j\sin(t/2)}\right|
 \le\frac{\pi}{j|t|}.
\]
Choosing a fixed $A>4\pi$ bounds every factor with $j\ge n/2$ by a constant below one. This proves the second bound, after harmless adjustment of constants.

The even cumulants are explicit:
\begin{equation}\label{mgp:eq:cumulants}
 \kappa_{2m}=\frac{B_{2m}}{2m}\sum_{j=1}^n(j^{2m}-1),\qquad
 \lambda_{2m}=\kappa_{2m}/\V^m=O_m(n^{1-m}).
\end{equation}
The convention is $B_2=1/6$, $B_4=-1/30$. For $u=\sigma_n t$,
\[
 \log\phi_n(u/\sigma_n)=-u^2/2+
 \sum_{m\ge2}\frac{(-1)^m\lambda_{2m}u^{2m}}{(2m)!}.
\]
For a prescribed order $L$ and finitely many derivatives, truncate the sum at a sufficiently large fixed $m$ and work first on $|u|\le n^\epsilon$, with $\epsilon>0$ sufficiently small. Taylor's theorem gives a remainder bounded by $n^{-L-1}$ times a fixed polynomial in $|u|$, multiplied by $e^{-c u^2}$. More explicitly, one may retain cumulants through $2(L+2)$, expand their rational functions of $1/n$ through $L+1$, and expand the exponential through weight $L+1$; the discarded terms have weight at least $L+2$. Decreasing $\epsilon$ makes the unexpanded remainder exponentially dominated by $e^{-c u^2}$. The region $n^\epsilon<|u|\le\sigma_n/n$ is smaller than any power by the first bound, and the rest by the second. Multiplication by any fixed power of $u$ is harmless. Integrating the resulting polynomials against the Gaussian proves the derivative statements. Positivity and the logarithmic derivative claim follow on each fixed interval because the limiting Gaussian is bounded away from zero there.
\end{proof}

Since $F_n$ is even, integrating its logarithmic second derivative twice gives, uniformly for $|x|\le B\sigma_n$ and $|x|\ge\delta$,
\begin{equation}\label{mgp:eq:profileleading}
 \log\frac{f_n(x)}{f_n(\delta)}
 =-\frac{x^2-\delta^2}{2\V}\{1+O_B(n^{-1})\}.
\end{equation}
This is a cancellation-respecting estimate: its error is proportional to $(x^2-\delta^2)/\V$. A relative local CLT with an undifferentiated $O(n^{-1})$ error would not suffice when raised to $r\asymp n^3$. Equivalently, the cancellation follows directly from
\[
 f_n(x)-f_n(\delta)=\frac1{2\pi}\int_{-\pi}^{\pi}
 \phi_n(t)\{\cos(xt)-\cos(\delta t)\}\dd t,
\]
where the difference contributes a factor $t^2$ for bounded $x$.

\section{Log concavity and the three regimes}
Each discrete uniform mass sequence is log concave and has interval support. Convolution preserves this property. For completeness, a sequence $b_k$ with these properties has every $2\times2$ minor of its Toeplitz matrix $(b_{i-j})$ nonnegative: the adjacent ratio $b_{k+1}/b_k$ is decreasing, and products of these ratios give the arbitrary minor inequality. The Toeplitz matrix of a convolution is the product of the two Toeplitz matrices. The $2\times2$ Cauchy--Binet formula preserves nonnegative minors; finite truncation suffices because the sequences have finite support. Thus the Mahonian sequence is log concave. Symmetry makes the one or two central lattice sites modes. The strict inequality $q_n<1$ proved below shows that, for sufficiently large $n$, there are no other modes.

Let $b_m=f_n(\delta+m)/f_n(\delta)$ for $m\ge0$ in the upper half of the support. The first ratio is $q_n=b_1$, and~\eqref{mgp:eq:profileleading} gives
\begin{equation}\label{mgp:eq:q}
 \log q_n=-\frac{1+2\delta}{2\V}\{1+O(n^{-1})\}.
\end{equation}
Log concavity implies $b_{m+1}/b_m\le q_n$, hence
\begin{equation}\label{mgp:eq:domination}
 0\le b_m\le q_n^m.
\end{equation}
For $\delta=1/2$, the ratio used here is the first ratio \emph{outside} the two equal modes: displacement $3/2$ divided by displacement $1/2$.

\begin{proof}[Proof of Theorem~\ref{mgp:thm:main}]
When $\rho_n$ remains in a compact positive interval,~\eqref{mgp:eq:profileleading} gives pointwise convergence of each summand, while~\eqref{mgp:eq:q}--\eqref{mgp:eq:domination} supply a geometric majorant $e^{-cm}$ independent of $n$. Dominated convergence proves the theta limit and its local uniformity.

If $\rho_n\to\infty$, symmetry and~\eqref{mgp:eq:domination} give
\begin{equation}\label{mgp:eq:firstshell}
 0\le T_{n,r}-(1+2\delta)-2q_n^r
 \le\frac{2q_n^{2r}}{1-q_n^r}.
\end{equation}
Equation~\eqref{mgp:eq:q} implies $q_n^r\to0$, proving the claimed limit and first-shell equivalent. This conclusion holds however fast $r$ grows because $q_n$ is kept exact.

Suppose finally that $r\to\infty$ and $\rho_n\to0$. On $|x|\le\sigma_n$,~\eqref{mgp:eq:profileleading} squeezes the summands between discrete Gaussian sums with parameters $\rho_n(1\pm C/n)$. Choose the last lattice point $M\le\sigma_n$. For large $n$, the same derivative estimate shows
\[
 f_n(M)/f_n(\delta)\le e^{-c},\qquad
 f_n(M)/f_n(M-1)\le e^{-c/\sigma_n}.
\]
Log concavity then bounds the sum outside this central interval by
\[
 O\left(\frac{e^{-cr}}{1-e^{-cr/\sigma_n}}\right)
 =O\bigl(e^{-cr}(1+\sigma_n/r)\bigr).
\]
Dividing by $\sigma_n/\sqrt r$ makes this $o(1)$; in this regime $\sqrt r/\sigma_n\to0$. The omitted Gaussian tails have the same negligible property. The elementary Gaussian lattice Riemann sum gives
\[
 \Theta_\delta(q)\sim\sqrt{2\pi/q}\qquad(q\downarrow0),
\]
uniformly for the two possible shifts. This proves the subcritical assertion. Every sequence of $\rho_n$ has a subsequence tending to zero, to a positive finite limit, or to infinity in the extended positive line; separating the two parity classes proves the full sequential statement.
\end{proof}

\section{Effective curvature and the explicit correction}
Define exact weighted Fourier moments
\[
 J_{h,n}=\frac1{\sqrt{2\pi}}
 \int_{-\pi\sigma_n}^{\pi\sigma_n}u^{2h}\phi_n(u/\sigma_n)\dd u.
\]
Then $F_n(0)=J_{0,n}/\sqrt{2\pi}$ and
\[
 a_n=\frac{J_{1,n}}{J_{0,n}},\qquad
 b_n=\frac{J_{2,n}}{J_{0,n}}-3a_n^2
\]
are respectively minus the second and the fourth derivatives at zero of $\log F_n$. Applying Lemma~\ref{mgp:lem:fourier} and Gaussian integration gives
\begin{align}
 a_n={}&1-\frac{27}{25n}+\frac{21087}{61250n^2}
 -\frac{19894023}{3062500n^3}
 +\frac{18889730796483}{907878125000n^4}+O(n^{-5}),\label{mgp:eq:a}\\
 b_n={}&-\frac{54}{25n}+O(n^{-2}).\label{mgp:eq:b}
\end{align}
The coefficient generator in the next section specifies the calculation without reliance on fitted values. In a fixed standardized neighborhood of zero,
\[
 \log F_n(y)-\log F_n(0)=-a_ny^2/2+b_ny^4/24+O(y^6).
\]
The remainder is uniform because the sixth logarithmic derivative is bounded. For $|x|\le A\log n$, it follows that
\begin{align}
 r\log\frac{f_n(x)}{f_n(\delta)}
 ={}&-\frac{\rho_n a_{[4]}(n)}2(x^2-\delta^2)
 -\frac{81\rho_n}{25n^4}(x^4-\delta^4)
 +O_K\bigl(n^{-5}(1+|x|^6)\bigr).\label{mgp:eq:exponent4}
\end{align}
Indeed, the quartic coefficient is $\rho_n b_n/(24\V)$, whose leading term is $-81\rho_n/(25n^4)$; the sixth-order remainder after multiplication by $r$ is $O_K(x^6\V^{-2})=O_K(n^{-6}x^6)$. Exponentiating~\eqref{mgp:eq:exponent4}, summing the polynomial remainder against a fixed geometric majorant, and using~\eqref{mgp:eq:domination} beyond $A\log n$ proves Theorem~\ref{mgp:thm:fourth}. The cutoff constant is chosen large enough that its geometric tail is $O_K(n^{-6})$. This summation uses polynomial weights inside the sum and therefore introduces no logarithmic loss.

\section{The arbitrary order generator}\label{mgp:sec:allorders}
The following is a formal power series prescription over the rationals. Put $z=1/n$ and
\[
 V(z)=\frac{1+3z/2-5z^2/2}{36z^3}.
\]
Evaluate the Faulhaber polynomials in~\eqref{mgp:eq:cumulants} at $n=1/z$ to obtain $\lambda_{2m}(z)$. With $U$ a formal standard Gaussian, set
\begin{align}
 A(u,z)&=\exp\left\{\sum_{m\ge2}
            \frac{(-1)^m\lambda_{2m}(z)u^{2m}}{(2m)!}\right\},\label{mgp:eq:formalA}\\
 J_h(z)&=\E\{U^{2h}A(U,z)\},\label{mgp:eq:formalJ}\\
 H(x,z)&=\sum_{h\ge0}\frac{(-1)^h x^{2h}}{(2h)!V(z)^h}
                      \frac{J_h(z)}{J_0(z)},\label{mgp:eq:formalH}\\
 P(x,z;\rho,\delta)&=\rho V(z)\{\log H(x,z)-\log H(\delta,z)\}
                    +\frac\rho2(x^2-\delta^2),\label{mgp:eq:formalP}\\
 A_j(x;\rho,\delta)&=[z^j]\exp P(x,z;\rho,\delta).
\end{align}
Here $\E U^{2k}=(2k-1)!!$, including $(-1)!!=1$. Every coefficient involves finitely many terms because $\lambda_{2m}=O(z^{m-1})$ and $V^{-1}=O(z^3)$. The constant term of $P$ is zero, and $A_0=1$. The coefficients $A_j$ are polynomials in $x,\delta,\rho$ with rational coefficients.

\begin{theorem}[Uniform finite expansions]\label{mgp:thm:allorders}
For each fixed integer $L\ge0$ and compact $K\subset(0,\infty)$,
\begin{equation}\label{mgp:eq:allorders}
 T_{n,r}=\sum_{j=0}^{L}\frac1{n^j}
 \sum_{x\in\Z+\delta}e^{-\rho_n(x^2-\delta^2)/2}
 A_j(x;\rho_n,\delta)+O_{K,L}(n^{-L-1}),
\end{equation}
uniformly for $\rho_n\in K$. Every inner sum is a finite linear combination of derivatives of $\Theta_\delta$.
\end{theorem}
\begin{proof}
Lemma~\ref{mgp:lem:fourier} proves that the exact moments $J_{h,n}$ have the expansions~\eqref{mgp:eq:formalJ}, to any prescribed order and for any fixed finite set of $h$. Taylor-expand $\log F_n(y)$ at zero through a sufficiently high even degree $2H$, with $3H\ge L+1$. All derivatives needed in the remainder are bounded uniformly for $|y|\le1$. Multiplication of that Taylor remainder by $r\asymp\V$ gives
\[
 O_K\bigl(\V^{-H}(1+|x|^{2H+2})\bigr)
 =O_K\bigl(n^{-L-1}(1+|x|^{2H+2})\bigr).
\]
For each retained Taylor coefficient, expand its weighted Fourier moments to enough inverse-$n$ orders. The quadratic term carries the factor $\V^{-1}$ before multiplication by $r$, so no factor $n^3$ is lost; higher Taylor terms are still smaller. This establishes, uniformly for $|x|\le A\log n$, the exponent expansion specified by~\eqref{mgp:eq:formalP}, with a remainder $O_{K,L}(n^{-L-1}(1+|x|^D))$ for a fixed $D$. The same statement after exponentiation follows because the correction to the leading Gaussian exponent tends to zero uniformly on this logarithmic window.

The actual summands on this window are bounded by a fixed geometric majorant by~\eqref{mgp:eq:domination}; equivalently, the leading Gaussian times its perturbation admits such a bound. Summing the polynomial error against this majorant gives $O_{K,L}(n^{-L-1})$, since $\sum_m(1+m^D)e^{-cm}<\infty$. Choose $A$ large enough that the actual tail is $O(n^{-L-2})$. The tails of every approximating polynomial times a Gaussian are even smaller. This proves~\eqref{mgp:eq:allorders}. Differentiation of~\eqref{mgp:eq:theta} generates all even powers of $x$ under the sum.
\end{proof}
This is an asymptotic expansion to arbitrary finite depth. No convergence of the infinite expansion, exponentially small remainder, or complete resurgent transseries is claimed.

\section{Removing the modal normalization}
The first coefficients from~\eqref{mgp:eq:formalJ} are
\[
 J_{0,n}=1-\frac{27}{100n}+\frac{1809}{196000n^2}
             -\frac{32054553}{19600000n^3}+O(n^{-4}),
\]
so
\begin{equation}\label{mgp:eq:logJ}
 \log J_{0,n}=-\frac{27}{100n}-\frac{6669}{245000n^2}
                    -\frac{20083941}{12250000n^3}+O(n^{-4}).
\end{equation}
At the lattice points, the exact representation is
\[
 \frac{M_n(\mu+x)}{n!}
 =\frac{J_{0,n}}{\sqrt{2\pi\V}}\frac{f_n(x)}{f_n(0)}.
\]
The same arguments as before, now normalizing at the Fourier interpolation's value at zero, give
\[
 \sum_{x\text{ in support}}\left(\frac{f_n(x)}{f_n(0)}\right)^r
 =Z_\delta(\rho_n)\{1+O_K(n^{-1})\}.
\]
This also follows directly from Theorem~\ref{mgp:thm:fourth} and
$r\log(f_n(\delta)/f_n(0))=-\rho_n\delta^2/2+O_K(n^{-1})$.
Finally, multiplying~\eqref{mgp:eq:logJ} by
\[
 r=\rho_n\V=\frac{\rho_n}{36}
              \left(n^3+\frac32n^2-\frac52n\right)
\]
produces
\[
 r\log J_{0,n}=-\rho_n\left(\frac{3n^2}{400}
                  +\frac{5883n}{490000}
                  +\frac{13963}{500000}\right)+O_K(n^{-1}),
\]
which proves Theorem~\ref{mgp:thm:unnorm}. For an error $O(n^{-L-1})$ in the final logarithm, retain~\eqref{mgp:eq:logJ} through $n^{-L-3}$; its next-order remainder multiplied by $r=O(n^3)$ has precisely that size. The lattice contribution is expanded by the same generator, omitting the subtraction at $\delta$.


\section{A fixed critical path and its inverse}
An inverse of the joint $(n,r)$ model is not determined until a path is specified. Fix $\rho>0$ and take the natural critical path $r_n=\rho\V$ exactly. Set
\[
 Y_n=\log S_n(\rho\V),\qquad
 X(y)=\left\{\frac{144y/\rho}{W_0((144y/\rho)e^{-4})}\right\}^{1/4},
\]
where $W_0$ is the positive principal Lambert branch for positive arguments. Then $X$ is the large positive solution of
\[
 \frac\rho{36}X^4(\log X-1)=y.
\]
The sequence $Y_n$ is strictly increasing for $n\ge2$: convolution gives $M_{n+1}(k)\ge M_n(k)$ on the old support, the support grows, all positive coefficients are at least one, and $r_{n+1}>r_n>0$.

\begin{theorem}[Fixed-path index inversion]\label{mgp:thm:inverse}
For exact index inputs $y=Y_n$,
\begin{equation}\label{mgp:eq:inverse}
 n=X(y)-\frac{\log X(y)+2\log6-3}{2(4\log X(y)-3)}
       +O_\rho(X(y)^{-1}).
\end{equation}
In particular, $n=X(y)-1/8+O_\rho(1/\log X(y))$. The expression on the right of~\eqref{mgp:eq:inverse} without its error term eventually recovers $n$ by nearest-integer rounding, but no certified numerical onset is asserted.
\end{theorem}
\begin{proof}
Stirling's formula and the exact variance polynomial give
\[
 \log\frac{n!}{\sqrt{2\pi\V}}
 =n\log n-n-\log n+\log6+O(n^{-1}).
\]
Using Theorem~\ref{mgp:thm:unnorm}, with $\rho$ fixed, therefore yields
\[
 Y_n=f(n)+g(n)+O_\rho(n^2\log n),
\]
where
\[
 f(x)=\frac\rho{36}x^4(\log x-1),\qquad
 g(x)=\frac\rho{72}x^3(\log x+2\log6-3).
\]
Since $f'(x)=\rho x^3(4\log x-3)/36$, the relation $f(X)=Y_n$ first gives $X/n\to1$, then $X-n=O(1)$. Taylor expansion now gives
\[
 n=X-\frac{g(X)}{f'(X)}+O_\rho(X^{-1}),
\]
which is~\eqref{mgp:eq:inverse}. The omitted $n^2\log n$ term and the quadratic Taylor correction both produce errors of order $X^{-1}$ after division by $f'(X)$.
\end{proof}

There is also a controlled arbitrary-order inverse without claiming a canonical real continuation of the original combinatorial sum. Write $\ell_j=[z^j]\log J_0(z)$ and let
\[
 B_{\delta,L}(x;\rho)=Z_\delta(\rho)+\sum_{j=1}^{L}B_{\delta,j}(\rho)x^{-j}
\]
be the order-$L$ lattice expansion normalized at $f_n(0)$. Its generator replaces~\eqref{mgp:eq:formalP} by $P(t,z)=\rho V(z)\log H(t,z)+\rho t^2/2$ and sums against $e^{-\rho t^2/2}$ on $\mathbb Z+\delta$. Define the explicit real model, for sufficiently large $x$, by
\begin{align*}
 \mathcal F_{\delta,L}(x;\rho)
 ={}&\rho V(x)\left\{\log\Gamma(x+1)-\frac12\log(2\pi V(x))
                    +\sum_{j=1}^{L+3}\ell_jx^{-j}\right\}\\
 &+\log B_{\delta,L}(x;\rho),\qquad
 V(x)=\frac{x(x-1)(2x+5)}{72}.
\end{align*}
Its lattice factor is eventually positive and its derivative is asymptotic to $f'(x)>0$. The forward theorem gives
\[
 Y_n=\mathcal F_{\delta,L}(n;\rho)+O_\rho(n^{-L-1})
\]
on the indicated parity class. Consequently the unique large real solution $x_{\delta,L}(y)$ of $\mathcal F_{\delta,L}(x;\rho)=y$ satisfies, at exact index inputs of that class,
\[
 x_{\delta,L}(Y_n)=n+O_\rho\left(\frac1{n^{L+4}\log n}\right).
\]
For $L=0$, this already uses only the three displayed $\ell_j$ in~\eqref{mgp:eq:logJ} and the familiar $Z_\delta$. A symbolic inverse in powers of $X^{-1}$ is obtained by inserting
$x=X+\sum_{m\ge0}X^{-m}R_m(\log X)$ and successively cancelling powers; the first coefficient is the correction in~\eqref{mgp:eq:inverse}. Each step is linear in the new coefficient, with nonzero divisor proportional to $4\log X-3$. The higher coefficients incorporate the fixed $\rho$ and parity-specific theta constants.

For arbitrary real threshold inputs $y$, the discrete answer is $\min\{n\ge2:Y_n\ge y\}$. A real approximation alone cannot decide rounding uniformly when $y$ is arbitrarily close to a threshold. Use the two parity models on their allowed residue classes, and check the neighboring integer candidates with the exact sequence, or with certified upper and lower error bounds. No unjustified ceiling rule is being substituted for a discrete inverse.

\begin{writenote}
The inversion mechanics of this section are instances of the repository's transseries volume
(\repo{Analysis/Transseries/docs/series-and-transseries/Transseries_And_Inversion/}). In the variable
$Z=\log X-1$ the defining equation $\frac\rho{36}X^4(\log X-1)=y$ reads
$4Z+\log Z=\log(36y/\rho)-4$, which is \texttt{p0:thm:lambert-core} with $a=4$, $b=1$ (branch $W_0$),
and gives $X(y)$ above. The correction $-g(X)/f'(X)$ of Theorem~\ref{mgp:thm:inverse} is the first-order term of
\texttt{p0:thm:perturbed-inversion} (with $F=f$ and $\varepsilon E=g$), which the symbolic inverse in
powers of $X^{-1}$ described above continues, and
the treatment of discrete thresholds in the last paragraph is the situation of
\texttt{p0:thm:staircase} (exact rounding for a monotone interpolation, and its separation
condition). No novelty is claimed for the inversion mechanics. Nothing in this report is
formalized; the generic staircase lemma is formalized as \texttt{Fabius.staircase\_separation} in
\repo{Analysis/FabiusFunction/Lean/FabiusFunction/StaircaseInversion.lean}, about an arbitrary
monotone interpolation, not about $S_n(r)$.
\end{writenote}

\section{Computational verification}
\begin{writenote}
``The accompanying programs'' are the five programs shipped byte-identical in the report's
\texttt{code/} directory; their recorded outputs are in \texttt{data/}, among them
\texttt{numerical\_checks.json} (the ``data file'' below: 30 critical-window cases, $n\le322$). The
delivered PDF is not shipped; \texttt{article.pdf} is a build of this text. The README gives a
rerun recipe on a scratch copy.
\end{writenote}
The accompanying programs use exact integer sliding-window recurrence for the full Mahonian rows, followed by 90-digit evaluation of normalized powers. No normal approximation is used to construct the reference values. Rational symbolic Gaussian moments independently generate the coefficients. The table shows $\rho_n=1$; the data file also covers $\rho_n=0.2$ and $5$. The final column is the relative error of Theorem~\ref{mgp:thm:fourth}'s displayed approximation.
\begin{center}
\begin{tabular}{rrrr}
\toprule
$n$&$\delta$&$T_{n,\V}$&relative fourth-order error\\
\midrule
20&$0.0$&$2.576804077525$&$8.565\times10^{-6}$\\
22&$0.5$&$2.894634886202$&$4.448\times10^{-6}$\\
80&$0.0$&$2.523667490870$&$9.898\times10^{-9}$\\
82&$0.5$&$2.854526220830$&$7.197\times10^{-9}$\\
320&$0.0$&$2.510864969401$&$1.011\times10^{-11}$\\
322&$0.5$&$2.843961378346$&$8.054\times10^{-12}$\\
\bottomrule
\end{tabular}
\end{center}
The raw theta approximation has a relative error of order $n^{-1}$. The fourth-order residuals scale as $n^{-5}$ in both parity classes. These checks support the algebra and parity conventions; the error bounds come from the proof, rather than numerical extrapolation.

\section{Literature boundary and interpretation}
For fixed $r$, $S_n(3)$ and $S_n(4)$ are OEIS A380274 and A380275. Wang~\cite{wang} already gives the fixed-real-power equivalent and three corrections; his Remark~2 expressly excludes varying $r$. The results here concern a moving parameter and do not resolve an unproved fixed-power leading term.

Local limits for Mahonian distributions have a substantial literature. Canfield, Janson and Zeilberger~\cite{cjz} prove local Gaussian estimates for q-multinomials; Section~4.1 obtains high-order central curvature for q-binomial coefficients. The use of Fourier expansion to recover microscopic curvature is therefore established methodology. Discrete Gaussian theta normalizers and their power identities are likewise standard~\cite{agostini,nielsen}. The contribution investigated here is the Mahonian-specific uniform growing-power approximation with exact modal normalization and explicit finite corrections.

\begin{writenote}
The research-report collection holds the same phenomenon for another array: its report
\texttt{a357825-theta-ballot-power-sums}, in the same category as this one, proves, for high powers of ballot numbers (OEIS A357824, A357825, A357871), a critical theta regime
with a varying phase (its Theorem \texttt{tbs:thm:critical}). Its research question Q6 (``A reusable
lattice-peak theorem for other path models'') asks for general hypotheses on a positive triangular
array under which high-power sums have a phase-dependent theta expansion. The present note is one
more instance of that phenomenon, for the Mahonian array, with only the two shifts
$\delta\in\{0,\tfrac12\}$; it does not prove the general theorem, and Q6 stays open. The two theta
normalizations differ: with $\Theta^{\mathrm{tbs}}_\tau(\alpha)=\sum_{j\in\Z}e^{-4\tau(j-\alpha)^2}$
there, $Z_\delta(q)=\Theta^{\mathrm{tbs}}_{q/8}(\delta)$ and
$\Theta_\delta(q)=e^{q\delta^2/2}\,\Theta^{\mathrm{tbs}}_{q/8}(\delta)$. The subscript of
$\Theta_\delta$ here is the lattice shift; the subscript of $\Theta^{\mathrm{tbs}}_\tau$ there is a
scale. Neither manuscript cites the other.
\end{writenote}

\begin{quote}\small\sloppy\textbf{[write, 2026-10-05]}\ The note above says
that question~Q6 of the ballot-sum report stays open. Since batch~98D that
report has a Part~II (its Section~12 onwards, labels \texttt{tbs:lpt:}),
which answers Q6 in a sufficient-hypothesis form: a general lattice-peak
transfer theorem (its Theorem~15.1, \texttt{tbs:lpt:thm:transfer}) under a
local quadratic law, a summable majorant and interior completeness. Its
Proposition~21.1 (\texttt{tbs:lpt:prop:mahonian}) shows that the critical
window of Theorem~\ref{mgp:thm:main} is an instance: \eqref{mgp:eq:profileleading}
gives the local law, and log-concavity of the Mahonian row (Section~3)
gives an exponential majorant, which its Theorem~15.4 shows is enough. That
proposition reorganizes the proof of Section~3 rather than giving an
independent one. Whether the rows also satisfy its Gaussian majorant is its
further question~F1. In its notation $Z_\delta(q)=\vartheta(q,\delta)$ and
$\Theta_\delta(q)=e^{q\delta^2/2}\vartheta(q,\delta)$.
\end{quote}

Vilenkin and D'yachkov~\cite{vd} treat entropy asymptotics for discrete and continuous i.i.d. sums. Only their primary abstract was available during this search, so its uniformity in R\'enyi order has not been verified. Bobkov and Marsiglietti~\cite{bm} treat fixed continuous R\'enyi orders, including a separate infinite-order endpoint. These references prevent a broad claim that high-order entropy limits themselves are new. A targeted search did not locate the specific theorem proved here, but this is not a priority certificate.

For a precise entropy interpretation, if $p_*=C_n/n!$, $H_\infty=-\log p_*$ and $H_r$ is the R\'enyi entropy of the inversion count, the exact identity is
\[
 (r-1)H_r-rH_\infty=-\log T_{n,r}.
\]
Thus the theta refinement remains order one in this scaled entropy difference. Its contribution to the unscaled $H_r$ is only of order $1/r$.

The result supplies a natural growing-tuple extension of two regular OEIS sequences. It is not itself a new asymptotic assertion for their fixed powers, and no beyond-all-orders expansion is asserted. The inverse section explicitly fixes a critical path and distinguishes the real approximation from discrete thresholds. The exact generating product and arbitrary finite coefficient procedure make further controlled extensions possible.

\begin{thebibliography}{9}
\bibitem{wang} X. Wang, \emph{Fixed-Power Sums of Mahonian Coefficients}, author manuscript, June 2026. \href{https://doi.org/10.5281/zenodo.20548010}{DOI 10.5281/zenodo.20548010}; \href{https://www.researchgate.net/publication/406002366_Fixed-Power_Sums_of_Mahonian_Coefficients}{author full text}.
\bibitem{cjz} E. R. Canfield, S. Janson and D. Zeilberger, \emph{The Mahonian probability distribution on words is asymptotically normal}, \href{https://arxiv.org/abs/0908.2089}{arXiv:0908.2089}, especially Theorem 4.5 and Section 4.1.
\bibitem{agostini} D. Agostini and C. Am\'endola, \emph{Discrete Gaussian distributions via theta functions}, \href{https://arxiv.org/abs/1801.02373}{arXiv:1801.02373}; \href{https://doi.org/10.1137/18M1164937}{DOI 10.1137/18M1164937}.
\bibitem{nielsen} F. Nielsen, \emph{The Kullback--Leibler Divergence Between Lattice Gaussian Distributions} (2022), \href{https://doi.org/10.1007/s41745-021-00279-5}{DOI 10.1007/s41745-021-00279-5}, Proposition 3.
\bibitem{vd} P. A. Vilenkin and A. G. D'yachkov, \emph{Asymptotics of the Shannon and Renyi Entropies for Sums of Independent Random Variables}, Problems of Information Transmission \textbf{34}(3) (1998), 17--31. \href{https://www.mathnet.ru/eng/ppi413}{Primary journal record}.
\bibitem{bm} S. G. Bobkov and A. Marsiglietti, \emph{Asymptotic Behavior of R\'enyi Entropy in the Central Limit Theorem}, High Dimensional Probability VIII (2019), Chapter 11. \href{https://arxiv.org/abs/1802.10212}{arXiv:1802.10212}.
\end{thebibliography}
\end{document}
